% Edited conversion of source pp. 35--36.
\ReportHeading
  {Алгоритм розв’язування осесиметричних задач термов’язкопластичності для тіл та оболонок із двокомпонентного функціонально-градієнтного матеріалу}
  {М.~О. Бабешко, В.~Г. Савченко}
  {Інститут механіки ім. С.~П. Тимошенка НАН України}
  {}

Розроблено алгоритм чисельного дослідження осесиметричного термов’язкопружнопластичного напружено-деформованого стану просторових тіл і тонких оболонок із двокомпонентного функціонально-градієнтного матеріалу (ФГМ) під час непружного деформування за неізотермічного навантаження.

Властивості матеріалу вважаються неперервними функціями координат. Закон зміни відносного об’ємного вмісту однієї зі складових задається наперед, а пружні й непружні характеристики визначаються об’ємними частками компонентів. Така постановка поширює підхід, використаний для пружних гнучких оболонок із ФГМ~\cite{fgmtvp:sklepus2023}, на крайові задачі термов’язкопластичності.

Сформульовано визначальні рівняння для ізотропного матеріалу, властивості якого неперервно змінюються за координатою, у межах теорії процесів деформування вздовж траєкторій малої кривизни. Повна деформація складається з пружної, теплової, пластичної та деформації повзучості. Тензори напружень і деформацій розкладаються на девіаторні та кульові частини; напрямні тензори девіаторів диференціалів непружних деформацій і напружень збігаються; пружні деформації пов’язані з напруженнями узагальненим законом Гука. У результаті зв’язок між компонентами напружень і деформацій має форму закону Гука з додатковими членами, які враховують теплову, пластичну та повзучу складові~\cite{fgmtvp:shevchenko2016,fgmtvp:monograph1987}.

У класичній теорії процесів малої кривизни функціональну залежність між другими інваріантами девіаторів напружень і деформацій вважають незалежною від виду напруженого стану та визначають за двома серіями експериментів: миттєвими діаграмами деформування за різних температур і кривими повзучості за різних напружень та температур. Для ФГМ у кожній точці тіла і за поточної температури відповідну діаграму деформування та сімейство кривих повзучості будують із діаграм компонентів та заданого закону зміни їх об’ємних часток. Ці локальні діаграми використовують для конкретизації інваріантної залежності на кожному кроці навантаження і в кожному наближенні.

Отже, визначальні рівняння ФГМ зберігають форму класичної теорії процесів малої кривизни, але відрізняються способом встановлення зв’язку між другими інваріантами девіаторів напружень і деформацій. Це дає змогу адаптувати раніше створені методи та програмні засоби аналізу непружного стану тіл і оболонок обертання. Адекватність моделі має оцінюватися шляхом порівняння розрахунків з експериментальними даними для конкретних функціонально-градієнтних матеріалів.

\begin{thebibliography}{9}
\bibitem{fgmtvp:sklepus2023}
S.~M. Sklepus,
``Method of solving geometrically nonlinear bending problems of functionally graded shallow shells of complex shape,''
\emph{International Applied Mechanics}, vol.~59, no.~6, pp.~725--733, 2023.

\bibitem{fgmtvp:shevchenko2016}
Yu.~N. Shevchenko and V.~G. Savchenko,
``Three-dimensional problems of thermoviscoplasticity: Focus on Ukrainian research (review),''
\emph{International Applied Mechanics}, vol.~52, no.~3, pp.~217--271, 2016.

\bibitem{fgmtvp:monograph1987}
Ю.~Н. Шевченко, В.~Г. Савченко,
\emph{Механика связанных полей в элементах конструкций. Т.~2: Термовязкопластичность}.
Киев: Наукова думка, 1987, 264~с.
\end{thebibliography}
